% !PRINTED 569
\weberpage{569}
%<<CONT
\weq{20}{
\begin{aligned}
(-1,\eta)^2 &= 17 \\
(i,\eta)^2 &= (-1,\eta)(-1+4i) \\
(\alpha,\eta)^2 &= -(i,\eta)(3+i\sqrt{8});
\end{aligned}
}

in addition one has $(1,\eta) = -1$,

\weq{21}{
\begin{aligned}
(i,\eta)(-i,\eta) &= 17 \\
(\alpha,\eta)(\alpha^{-1},\eta) &= 17 \\
(\alpha^3,\eta)(\alpha^{-3},\eta) &= 17
\end{aligned}
}

and by (9)

\weq{22}{
(\alpha,\eta)(\alpha^3,\eta) = (-1,\eta)(3+i\sqrt{8}),
}

whereby $(1,\eta)$, $(-1,\eta)$, $(\pm i,\eta)$, $(\alpha,\eta)$, $(\alpha^{-1},\eta)$, $(\alpha^{-3},\eta)$ are
determined by square roots.
% [source-original mathematical omission, verified 2026-08-29] Formula (22), together
% with the already determined (alpha,eta), also determines (alpha^3,eta), but the printed
% enumeration omits that quantity while retaining its reciprocal partner (alpha^-3,eta).

With the determination that we have made concerning $r$ and $\alpha$,
it follows easily from a consideration of the values of the cosines
(16) that $(-1,\eta)$ is positive, and hence equal to the positive square root
$\sqrt{17}$, and further that $(i,\eta)$ and $(\alpha,\eta)$ have positive real
parts, whereby these quantities too are completely determined.
We have therefore

\[
\begin{aligned}
(-1,\eta) &= \sqrt{17} \\
(i,\eta) &= \sqrt[4]{17}\left\{\sqrt{\frac{\sqrt{17}-1}{2}} + i\sqrt{\frac{\sqrt{17}+1}{2}}\right\},
\end{aligned}
\]

if all the roots are taken positive; and by the last
equation (20), $(\alpha,\eta)$ can be represented, by a formula which is,
to be sure, somewhat long, through square roots of real quantities.
We shall not enter into the geometric constructions that
attach themselves here\footnote{Cf. on this v. \gesperrt{Staudt}, Construction des regulären Siebzehnecks. Crelle's Journ. Bd. 24, 1842.}.

The determination of signs employed is obtained simply
if, following (16), one sets

\[
\frac{1}{2}(-1,\eta) = \cos\frac{2\pi}{17} - \cos\frac{\pi}{17} + \cos\frac{8\pi}{17} + \cos\frac{4\pi}{17}
- \cos\frac{6\pi}{17} + \cos\frac{3\pi}{17} + \cos\frac{7\pi}{17} + \cos\frac{5\pi}{17}.
\]

But by a well-known trigonometric formula,

\[
\begin{aligned}
\cos\frac{2\pi}{17} - \cos\frac{\pi}{17} &= -2\sin\frac{\pi}{34}\sin\frac{3\pi}{34} \\
-\cos\frac{6\pi}{17} + \cos\frac{5\pi}{17} &= 2\sin\frac{\pi}{34}\sin\frac{11\pi}{34},
\end{aligned}
\]
%CONT>>
